\section{Finite-information prediction: Gaussian reduction and minimax
certificates}\label{gaussian_decision-finite-information-prediction-gaussian-reduction-and-minimax-certificates}

Extension of L.1--L.8 for the known scalar entropy-tilt family. Prepared
29 October 2026.

\subsection{1. Observation experiment and
notation}\label{gaussian_decision-observation-experiment-and-notation}

Let the calibration observations have density f\_theta=f0(1+theta h),
with bounded h, E0 h=0 and E0 h\^{}2=1. In the explicit entropy
construction, f0 is a centered Gaussian and h(z)=U(0,z)/sigma\_g. For
each n the parameter is u in the fixed compact interval {[}-A0,A0{]},
theta=u/sqrt(n), and the calibration experiment is

\[
\mathcal E_n=\{P_{n,u}=Q_{u/\sqrt n}^{\otimes M}:|u|\le A_0\}.
\]

Assume M/n-\textgreater rho in (0,infinity). The true target is the
probability pi\_n(Q\_\{u/sqrt(n)\}) of right risk exit by a future,
independent entropy-SGD trajectory. The supplied uniform exit lemma
gives

\[
r_n:=\sup_{|u|\le A_0}|\pi_n(Q_{u/\sqrt n})-\Phi(cu)|\to0.
\]

Here c=sqrt((2+eta lambda)/(eta lambda)) is fixed. Define the Gaussian
experiment

\[
\mathcal G_\rho=\{Y=u+\rho^{-1/2}Z:|u|\le A_0\},\qquad Z\sim N(0,1).
\]

Throughout this extension TV means sup\_B \textbar P(B)-Q(B)\textbar, or
half the L1 distance of densities. Nussbaum's reference uses the full L1
convention, so constants must be translated when comparing definitions.
We use deficiency as the infimum over parameter-independent Markov
kernels of the supremum TV approximation error, and Delta as the maximum
of the two deficiencies.

Reference: Michael Nussbaum, \emph{Asymptotic Equivalence of Statistical
Experiments}, especially the definitions of deficiency and its
decision-theoretic interpretation, pp.~4--6:
https://pi.math.cornell.edu/\textasciitilde nussbaum/papers/04-1-transl.pdf
. The proof below is specialized to the present bounded tilt family; the
reference is not being cited as proving this model-specific result.

\subsection{2. Theorem L.9: equivalence of the entire local
experiment}\label{gaussian_decision-theorem-l.9-equivalence-of-the-entire-local-experiment}

For the explicit family and every fixed A0\textless infinity,

\[
\boxed{\Delta(\mathcal E_n,\mathcal G_\rho)\to0.}
\tag{G.1}
\]

The assertion supplies both directions of risk transfer, not merely a
CLT for the empirical moment.

\textbf{Proof.} Write P\_\{n,0\}=Q0\^{}\{otimes M\}, rho\_n=M/n, and

\[
S_n=M^{-1/2}\sum_{j=1}^M h(Z_j),\qquad
V_n=M^{-1}\sum_{j=1}^M h(Z_j)^2.
\]

Let F\_M be the CDF of S\_n under P\_\{n,0\}. In the present family h is
a nonconstant real-analytic function, so the preimage of any value is a
Lebesgue-null set. Since f0 has a density, h(Z) has no atoms; the sum
also has no atoms. Therefore F\_M is continuous. Define

\[
T_n=\Phi^{-1}(F_M(S_n)).
\]

This parameter-independent statistic has the exact N(0,1) law under
P\_\{n,0\}; events where the argument is zero or one have probability
zero. The baseline law and the transform are known because the model
family is known.

The CLT and continuity of Phi give
sup\_x\textbar F\_M(x)-Phi(x)\textbar-\textgreater0. On a compact
interval for S\_n the inverse normal CDF is uniformly continuous after
restricting to a slightly larger compact probability interval. Tightness
of S\_n then proves

\[
S_n-T_n\longrightarrow0\quad\text{in }P_{n,0}\text{-probability}.
\tag{G.2}
\]

This use of a CLT is only one part of the proof. The essential next step
controls the full likelihood ratios.

Let Lambda\_n(u)=dP\_\{n,u\}/dP\_\{n,0\}. Boundedness of h and
compactness of u allow a uniform Taylor expansion, valid for all
sufficiently large n:

\[
\log\Lambda_n(u)
=u\sqrt{\rho_n}S_n-\tfrac12u^2\rho_n V_n+R_n(u),
\qquad
\sup_{|u|\le A_0}|R_n(u)|\le C M n^{-3/2}.
\tag{G.3}
\]

Under P\_\{n,0\}, V\_n-\textgreater1 in probability, and
rho\_n-\textgreater rho. Equations (G.2)--(G.3) imply

\[
\sup_{|u|\le A_0}
\left|\log\Lambda_n(u)-\left(\sqrt\rho\,uT_n-\tfrac12\rho u^2\right)\right|
\longrightarrow0
\tag{G.4}
\]

in probability. Define g\_u(t)=exp(sqrt(rho)u t-rho u\^{}2/2). We now
strengthen convergence of log likelihoods to uniform L1 convergence of
likelihoods.

For any fixed p\textgreater1, log(1+x)\textless=x gives

\[
\sup_{|u|\le A_0}\Lambda_n(u)
\le\exp\left(A_0\left|\sum_jh(Z_j)\right|/\sqrt n\right).
\]

The centered bounded summands have a sub-Gaussian exponential-moment
bound. Since M/n is bounded, the pth moments of the last display are
uniformly bounded. Also sup\_u
g\_u(T\_n)\textless=exp(sqrt(rho)A0\textbar T\_n\textbar), whose pth
moment is bounded because T\_n is exactly standard normal. Equation
(G.4), tightness, and these bounds give uniform integrability and hence

\[
\varepsilon_n:=\tfrac12\mathbb E_{n,0}
\sup_{|u|\le A_0}|\Lambda_n(u)-g_u(T_n)|\longrightarrow0.
\tag{G.5}
\]

Define an auxiliary experiment on the original sample space by

\[
d\widetilde P_{n,u}=g_u(T_n)\,dP_{n,0}.
\]

These are probability measures, because E\_\{n,0\}g\_u(T\_n)=1 exactly.
Equation (G.5) bounds sup\_u TV(P\_\{n,u\},Ptilde\_\{n,u\}) by
epsilon\_n.

Under Ptilde\_\{n,u\}, T\_n has law N(sqrt(rho)u,1). Conversely, use the
regular conditional distribution P\_\{n,0\}(d z\textbar T\_n=t) to turn
such a normal observation back into a full calibration sample. This
kernel is independent of u; tilting by a function of T\_n leaves the
conditional law unchanged. Regular conditional distributions exist on
the Euclidean sample spaces, and arbitrary versions on baseline-null
sets do not matter because all normal shifts are mutually absolutely
continuous. Thus the auxiliary experiment is exactly equivalent to
observing T\textasciitilde N(sqrt(rho)u,1), or Y=T/sqrt(rho).

Applying the statistic T\_n in one direction and this baseline
conditional kernel in the other gives both deficiencies bounded by
epsilon\_n.~This proves (G.1). \(\square\)

The proof supplies asymptotic equivalence on a fixed compact local
parameter interval. It does not give a calibrated numerical value of
epsilon\_n, nor global equivalence over an expanding parameter range.
The conditional kernel is an existence construction; implementing it is
not required to use a decision rule in the forward direction.

\subsection{3. Corollary L.10: the finite-ratio minimax
limit}\label{gaussian_decision-corollary-l.10-the-finite-ratio-minimax-limit}

Let

\[
V(A_0,c,\rho)=\inf_\delta\sup_{|u|\le A_0}
 \mathbb E_u|\delta(Y)-\Phi(cu)|.
\]

Then the original minimax risk for predicting the true future SGD
direction probability satisfies

\[
\boxed{\mathfrak R_{n,M}(A_0)\longrightarrow V(A_0,c,\rho).}
\tag{G.6}
\]

\textbf{Proof.} Project all decision values to {[}0,1{]}, which cannot
increase loss. Absolute loss is then bounded by one, so transporting any
randomized rule through either Markov kernel changes each parameterwise
risk by at most epsilon\_n.~Replacing the target pi\_n by Phi(cu)
changes each loss by at most r\_n.~Taking supremum risks and then infima
in both directions gives

\[
|\mathfrak R_{n,M}(A_0)-V(A_0,c,\rho)|\le\varepsilon_n+r_n.
\]

This reasoning does not require the minimax infimum to be attained:
apply it to arbitrarily near-optimal rules. Under absolute loss,
averaging any external decision randomization conditionally on the
observation cannot worsen risk, by convexity. Thus restricting to
deterministic rules leaves the values unchanged. \(\square\)

Standardize with

\[
t=cu,\qquad K=cA_0,\qquad X=cY=t+rZ,\qquad r=c/\sqrt\rho.
\]

The limiting decision problem depends only on (K,r):

\[
V_G(K,r)=\inf_d\sup_{|t|\le K}\mathbb E_t|d(X)-\Phi(t)|.
\tag{G.7}
\]

Actions may be restricted further to {[}Phi(-K),Phi(K){]}. Reflection
and averaging permit a symmetric rule d(-x)=1-d(x) without increasing
worst-case risk, although the certificates below do not require a proof
that the candidate optimizer is minimax.

\subsection{4. Certified numerical brackets without claiming optimizer
optimality}\label{gaussian_decision-certified-numerical-brackets-without-claiming-optimizer-optimality}

\subsubsection{4.1. A prior produces a valid lower
bound}\label{gaussian_decision-a-prior-produces-a-valid-lower-bound}

Choose any finite prior with ordered support
t\_1\textless\ldots\textless t\_m in {[}-K,K{]} and positive weights
w\_i. Its posterior median under absolute loss is a Bayes rule taking
the values Phi(t\_i). The boundary between actions k and k+1 solves

\[
\sum_{i\le k}w_i e^{x t_i/r^2-t_i^2/(2r^2)}
=\sum_{i>k}w_i e^{x t_i/r^2-t_i^2/(2r^2)}.
\tag{G.8}
\]

The ratio of the left sum to the right sum is strictly decreasing: its
log derivative is the difference of two tilted means divided by r\^{}2,
and every left support point is smaller than every right support point.
Its limits are infinity and zero. Thus every crossing is unique; the
ordered crossings define the posterior-median rule.

For any step rule with thresholds b\_j and action values a\_j, its exact
risk is

\[
R_d(t)=\sum_j |a_j-\Phi(t)|
\left[\Phi\left(\frac{b_{j+1}-t}{r}\right)
      -\Phi\left(\frac{b_j-t}{r}\right)\right],
\tag{G.9}
\]

including b\_0=-infinity and the last boundary +infinity. Therefore the
full unbounded observation domain is integrated through CDF differences,
with no omitted integration tail.

The code optimizes a finite prior using floating-point quadrature only
to find a candidate. It then rounds support points, prior numerators and
rule thresholds to explicit finite decimals. These rounded quantities
define the actual certified prior and rule.

For each true crossing, directed interval arithmetic verifies opposite
signs of (G.8) at endpoints l\_j and u\_j, with the candidate threshold
inside the bracket. The candidate rule and exact Bayes rule can differ
only on the union of these strips. Gaussian mixture density is at most
1/(r sqrt(2 pi)); loss differences are at most one. Hence

\[
\text{Bayes risk}\ge
\sum_i w_i R_d(t_i)-\frac{\sum_j(u_j-l_j)}{r\sqrt{2\pi}}.
\tag{G.10}
\]

A lower enclosure for the right side is a valid minimax lower bound.
This does not assume that the candidate prior is least favorable, that
the search objective was computed exactly, or that the optimizer
converged.

\subsubsection{4.2. Continuous-window upper
bound}\label{gaussian_decision-continuous-window-upper-bound}

The candidate step rule itself gives an upper bound once
sup\_\{\textbar t\textbar\textless=K\} R\_d(t) is bounded. Its action
values are exactly Phi(t\_i), so the only absolute-loss kinks are at the
known support points t\_i. Between successive kinks, differentiation
under the Gaussian integral gives

\[
|R_d''(t)|\le C_2:=1+2/r+2/r^2.
\tag{G.11}
\]

Indeed the target's first and second derivatives are bounded in absolute
value by one; the sum of absolute bin-probability derivatives is bounded
by E\textbar Z\textbar/r\textless=1/r, and that of second derivatives by
E\textbar Z\textsuperscript{2-1\textbar/r}2\textless=2/r\^{}2. All
action losses are bounded by one. The bound is deliberately
conservative.

For a partition interval {[}a,b{]} with no kink in its interior, linear
interpolation and (G.11) imply

\[
\sup_{t\in[a,b]}R_d(t)
\le\max\{R_d(a),R_d(b)\}+C_2(b-a)^2/8.
\tag{G.12}
\]

The partition includes every support point and both endpoints. Adaptive
refinement produces an upper bound over the entire continuous window,
not just the evaluated grid.

\subsubsection{4.3. Arithmetic and probability
enclosures}\label{gaussian_decision-arithmetic-and-probability-enclosures}

\texttt{certify.py} uses Python Decimal with 90-digit precision and
outward rounding for interval operations. Correctly rounded Decimal
exponentials and square roots are expanded by a representable value on
each side. Pi is enclosed by Machin's identity and alternating
arctangent remainders. Normal CDF values are enclosed by integrating the
Taylor expansion of exp(-x\^{}2/2), including a rigorous next-term
remainder once the terms decrease. For \textbar x\textbar\textgreater=12
a Mills-ratio tail bound is used instead. Both negative and positive
tails are retained in their intervals.

The optimizer's normal CDF, integration grid, stopping tolerance and
floating-point objective do not enter the validity of the final bounds.
The certificate relies on the stated mathematical bounds, the
correctness of the supplied implementation, and the documented Decimal
arithmetic; it is not a proof-assistant verification of the runtime.

The search gap and the certified gap are different quantities. The final
upper-minus-lower interval accounts for an imperfect prior/rule pair as
well as root, probability-evaluation and continuous-window errors.

\subsection{5. Using the computed Gaussian rule on calibration
data}\label{gaussian_decision-using-the-computed-gaussian-rule-on-calibration-data}

For an arbitrary measurable Gaussian rule, the direct risk-transfer
construction applies it to c T\_n/sqrt(rho). This uses the known
baseline CDF F\_M, not the unknown environment parameter.

For each of the fixed finite step rules produced here, a simpler
implementation is also asymptotically justified: apply it to Xhat=c
sqrt(n) hat\_theta, the raw empirical moment coordinate. Under the
baseline, Xhat-c T\_n/sqrt(rho) tends to zero in probability. The
likelihood-ratio moment bounds used in (G.5) transfer this to the
compact local family uniformly. The Gaussian statistic has a continuous
density uniformly bounded on the local window, so the probability of
lying near one of the finitely many thresholds tends to zero with the
threshold neighborhood. Therefore the two implementations have uniformly
vanishing risk difference.

This argument is for a fixed finite rule. Refining the rule with n so
that the number of thresholds diverges requires additional control. No
finite-n risk certificate for the SGD experiment follows merely from the
Gaussian numerical certificate.

\subsection{6. Reproducible files and validation
boundaries}\label{gaussian_decision-reproducible-files-and-validation-boundaries}

\begin{itemize}
\tightlist
\item
  \texttt{solve.py}: exploratory prior optimization for six (K,r)
  settings. A fixed support grid and observation quadrature generate
  candidates; neither is treated as a minimax proof.
\item
  \texttt{certify.py}: independent outward interval evaluation of a
  prior lower bound and a continuous-parameter rule upper bound.
\item
  \texttt{candidates.json}: floating-point search outputs.
\item
  \texttt{certificates.json}: exact finite-decimal prior/rule
  specifications, crossing brackets, and outward risk bounds.
\item
  \texttt{sgd\_pilot.py}: independent calibration and actual entropy-SGD
  comparison for a fixed selected rule.
\item
  \texttt{sgd\_results.json} and \texttt{pilot\_*.npz}: pilot summary,
  exit directions/times, and all four forecast arrays.
\end{itemize}

Run with Python, NumPy and SciPy; Decimal is in the standard library.
The checks in this run used NumPy 2.3.5 and SciPy 1.17.0.
\texttt{OPENBLAS\_NUM\_THREADS=1\ OMP\_NUM\_THREADS=1} avoids
unnecessary threaded overhead in the small optimizations.

The formal progress is the experiment-level equivalence and the
resulting minimax convergence, plus numerical brackets for the Gaussian
limit. The pilot tests their relevance at finite n; its calibration and
Monte Carlo errors are not part of those proofs.

\subsection{7. Results}\label{gaussian_decision-results}

Numerical tables and the actual-SGD pilot summary are appended below
after validation.

\subsubsection{7.1. Gaussian minimax
intervals}\label{gaussian_decision-gaussian-minimax-intervals}

Bounds below are rounded outward to six decimals. Raw-moment and
constant risks have the analytic formulas arctan(r)/pi and Phi(K)-.5.
The clipped-rule column is floating-point quadrature plus numerical
maximization and is marked approximate; it is not needed for either
minimax certificate.

{\def\LTcaptype{table} % do not increment counter
\begin{longtable}[]{@{}
  >{\raggedleft\arraybackslash}p{(\linewidth - 10\tabcolsep) * \real{0.1667}}
  >{\raggedleft\arraybackslash}p{(\linewidth - 10\tabcolsep) * \real{0.1667}}
  >{\raggedleft\arraybackslash}p{(\linewidth - 10\tabcolsep) * \real{0.1667}}
  >{\raggedleft\arraybackslash}p{(\linewidth - 10\tabcolsep) * \real{0.1667}}
  >{\raggedleft\arraybackslash}p{(\linewidth - 10\tabcolsep) * \real{0.1667}}
  >{\raggedleft\arraybackslash}p{(\linewidth - 10\tabcolsep) * \real{0.1667}}@{}}
\toprule\noalign{}
\begin{minipage}[b]{\linewidth}\raggedleft
K
\end{minipage} & \begin{minipage}[b]{\linewidth}\raggedleft
r
\end{minipage} & \begin{minipage}[b]{\linewidth}\raggedleft
Certified minimax interval
\end{minipage} & \begin{minipage}[b]{\linewidth}\raggedleft
Constant worst risk
\end{minipage} & \begin{minipage}[b]{\linewidth}\raggedleft
Raw moment worst risk
\end{minipage} & \begin{minipage}[b]{\linewidth}\raggedleft
Clipped worst risk (approx.)
\end{minipage} \\
\midrule\noalign{}
\endhead
\bottomrule\noalign{}
\endlastfoot
1 & 2 & {[}0.226230, 0.226386{]} & 0.341345 & 0.352416 & 0.279891 \\
1 & 1 & {[}0.164401, 0.164587{]} & 0.341345 & 0.250000 & 0.224829 \\
1 & 0.5 & {[}0.107333, 0.107539{]} & 0.341345 & 0.147584 & 0.145865 \\
2 & 2 & {[}0.244686, 0.244931{]} & 0.477250 & 0.352416 & 0.346916 \\
2 & 1 & {[}0.169889, 0.170268{]} & 0.477250 & 0.250000 & 0.249482 \\
2 & 0.5 & {[}0.108534, 0.109092{]} & 0.477250 & 0.147584 & 0.147583 \\
\end{longtable}
}

For K=1,r=1, the certified optimized rule has worst-case risk at most
0.164587, versus exactly 0.25 for the raw moment rule: at least a
34.16\% reduction in Gaussian worst-case MAE. This is not a claim of
pointwise dominance, and is not a finite-n guarantee for entropy SGD.

\subsubsection{7.2. Actual entropy-SGD
pilot}\label{gaussian_decision-actual-entropy-sgd-pilot}

The pilot fixes w=2, nu=1, eta=.1, d=.5, b=.25, kappa=.5,
epsilon=0.0060745243, and x0=0. Numerical integration gives
lambda=0.2421101333, sigma\_g\^{}2=0.1078103527, and c=9.1436882021. The
selected Gaussian rule has K=r=1; its prior and thresholds were
determined entirely in the Gaussian experiment before viewing SGD
results. Each n uses M=round(c\^{}2 n), and t in \{-1,-.5,0,.5,1\}, with
theta=t/(c sqrt(n)).

Each cell uses 1,024 independent true-SGD reference paths and 256
independent calibration repetitions; the four rules share calibration
samples within a repetition. Reference and calibration streams are
independent. Rejection sampling draws from the actual tilted continuous
density. Updates use sampled nonlinear entropy gradients, not a Gaussian
autoregression or interpolated population field. The 1,500-step cap
censored no paths in any cell. Raw exit directions/times and forecasts
are retained.

{\def\LTcaptype{table} % do not increment counter
\begin{longtable}[]{@{}
  >{\raggedleft\arraybackslash}p{(\linewidth - 12\tabcolsep) * \real{0.1429}}
  >{\raggedleft\arraybackslash}p{(\linewidth - 12\tabcolsep) * \real{0.1429}}
  >{\raggedleft\arraybackslash}p{(\linewidth - 12\tabcolsep) * \real{0.1429}}
  >{\raggedleft\arraybackslash}p{(\linewidth - 12\tabcolsep) * \real{0.1429}}
  >{\raggedleft\arraybackslash}p{(\linewidth - 12\tabcolsep) * \real{0.1429}}
  >{\raggedleft\arraybackslash}p{(\linewidth - 12\tabcolsep) * \real{0.1429}}
  >{\raggedleft\arraybackslash}p{(\linewidth - 12\tabcolsep) * \real{0.1429}}@{}}
\toprule\noalign{}
\begin{minipage}[b]{\linewidth}\raggedleft
n
\end{minipage} & \begin{minipage}[b]{\linewidth}\raggedleft
M
\end{minipage} & \begin{minipage}[b]{\linewidth}\raggedleft
Largest cell MAE: constant
\end{minipage} & \begin{minipage}[b]{\linewidth}\raggedleft
Raw moment
\end{minipage} & \begin{minipage}[b]{\linewidth}\raggedleft
Clipped
\end{minipage} & \begin{minipage}[b]{\linewidth}\raggedleft
Optimized
\end{minipage} & \begin{minipage}[b]{\linewidth}\raggedleft
Maximum exit step
\end{minipage} \\
\midrule\noalign{}
\endhead
\bottomrule\noalign{}
\endlastfoot
128 & 10702 & 0.338867 & 0.253404 & 0.224829 & 0.167207 & 685 \\
512 & 42807 & 0.352539 & 0.267974 & 0.238485 & 0.178871 & 534 \\
\end{longtable}
}

These are maxima of empirical means over five parameter values, not
continuous-window upper bounds. Calibration sampling uncertainty and
reference-probability uncertainty remain. In particular, the exact
target at t=0 is .5 by reflection symmetry; the table consistently uses
the independent empirical reference in all cells, so the constant rule
has a small nonzero reported center error. No claim is made that these
runs verify the narrow Gaussian certificate.

{\def\LTcaptype{table} % do not increment counter
\begin{longtable}[]{@{}
  >{\raggedleft\arraybackslash}p{(\linewidth - 12\tabcolsep) * \real{0.1429}}
  >{\raggedleft\arraybackslash}p{(\linewidth - 12\tabcolsep) * \real{0.1429}}
  >{\raggedleft\arraybackslash}p{(\linewidth - 12\tabcolsep) * \real{0.1429}}
  >{\raggedleft\arraybackslash}p{(\linewidth - 12\tabcolsep) * \real{0.1429}}
  >{\raggedleft\arraybackslash}p{(\linewidth - 12\tabcolsep) * \real{0.1429}}
  >{\raggedleft\arraybackslash}p{(\linewidth - 12\tabcolsep) * \real{0.1429}}
  >{\raggedleft\arraybackslash}p{(\linewidth - 12\tabcolsep) * \real{0.1429}}@{}}
\toprule\noalign{}
\begin{minipage}[b]{\linewidth}\raggedleft
n
\end{minipage} & \begin{minipage}[b]{\linewidth}\raggedleft
t
\end{minipage} & \begin{minipage}[b]{\linewidth}\raggedleft
Observed right probability
\end{minipage} & \begin{minipage}[b]{\linewidth}\raggedleft
Gaussian Phi(t)
\end{minipage} & \begin{minipage}[b]{\linewidth}\raggedleft
Optimized MAE
\end{minipage} & \begin{minipage}[b]{\linewidth}\raggedleft
Raw minus optimized MAE
\end{minipage} & \begin{minipage}[b]{\linewidth}\raggedleft
Paired calibration SE
\end{minipage} \\
\midrule\noalign{}
\endhead
\bottomrule\noalign{}
\endlastfoot
128 & -1 & 0.176758 & 0.158655 & 0.155629 & 0.027099 & 0.006051 \\
128 & -0.5 & 0.303711 & 0.308538 & 0.161231 & 0.069069 & 0.004554 \\
128 & 0 & 0.498047 & 0.500000 & 0.167207 & 0.086197 & 0.002743 \\
128 & 0.5 & 0.687500 & 0.691462 & 0.166523 & 0.066868 & 0.004261 \\
128 & 1 & 0.838867 & 0.841345 & 0.163469 & 0.028516 & 0.005741 \\
512 & -1 & 0.147461 & 0.158655 & 0.163874 & 0.007559 & 0.005150 \\
512 & -0.5 & 0.335938 & 0.308538 & 0.155958 & 0.072129 & 0.004199 \\
512 & 0 & 0.494141 & 0.500000 & 0.176939 & 0.091036 & 0.002541 \\
512 & 0.5 & 0.698242 & 0.691462 & 0.148184 & 0.059938 & 0.004745 \\
512 & 1 & 0.851562 & 0.841345 & 0.178871 & 0.009024 & 0.005581 \\
\end{longtable}
}

Paired standard errors condition on the common empirical reference
probability; they do not include its uncertainty. The reference Wilson
intervals, supplied in the JSON, are descriptive per-cell intervals
rather than simultaneous certification. With 1,024 paths, center-cell
reference uncertainty is about .03 at the usual 95\% scale, much larger
than the Gaussian minimax bracket widths. A two-rule loss difference can
change by up to twice the reference error.

The geometry diagnostic evaluates v0 prime on a grid of spacing .005 in
{[}-.5,.5{]}, giving a numerical minimum .19440797. The analytic bound
\textbar v0 double-prime\textbar\textless=6 gives an
interpolation-corrected lower value .17940797, above lambda/2=.12105507.
An additional conservative tilt correction gives .16468778 over the
pilot parameter range, above lambda/4=.06052753. The derivative bound
follows from \textbar a prime\textbar\textless=a, \textbar a
double-prime\textbar\textless=a and
\textbar t\textbar a(t)\textless=1/4. The quadrature itself was not
enclosed with directed intervals, so this is a diagnostic geometry
check, separate from the certified Gaussian decision calculation.

At the window endpoints, clipping can beat the optimized rule, while the
optimized rule substantially reduces the central and worst-grid errors.
That pattern is consistent with minimax balancing rather than pointwise
domination. The next experimental priority is more calibration
repetitions and reference paths at fixed predeclared cells if precise
finite-n improvements are required; there is no reason yet to add
stopping strategies.
